\section*{Capítulo \ref{ch:g12:seq} --- Sequências}

\begin{solution}{exo:g12:seq:1}
Seja $P(n)$ a afirmação
$\sum_{k=1}^{n} k^2 = \frac{n(n+1)(2n+1)}{6}$.
\emph{Caso base:} para $n = 0$ os dois lados valem $0$ (soma vazia).
\emph{Passo indutivo:} suponha $P(n)$. Então
\[
\sum_{k=1}^{n+1} k^2 = \frac{n(n+1)(2n+1)}{6} + (n+1)^2
= \frac{(n+1)\bigl(n(2n+1) + 6(n+1)\bigr)}{6}
= \frac{(n+1)(2n^2 + 7n + 6)}{6}.
\]
Como $2n^2 + 7n + 6 = (n+2)(2n+3)$, isso é
$\frac{(n+1)(n+2)(2(n+1)+1)}{6}$, que é $P(n+1)$. Por indução, $P(n)$
vale para todo $n$.
\end{solution}

\begin{solution}{exo:g12:seq:2}
$a_{n+1} - a_n = \frac{n+2}{n+1} - \frac{n+1}{n}
= \frac{n(n+2) - (n+1)^2}{n(n+1)} = \frac{-1}{n(n+1)} < 0$: $(a_n)$ é
estritamente \omterm{def:g10:functions:variations}{decrescente}.

$(b_n)$ tem termos positivos e
$\frac{b_{n+1}}{b_n} = \frac{2^{n+1}}{n+1}\cdot\frac{n}{2^n} =
\frac{2n}{n+1} \geq 1 \iff 2n \geq n+1 \iff n \geq 1$: $(b_n)$ é \omterm{def:g12:seq:monotonic}{crescente}
(estritamente para $n \geq 2$).

$c_{n+1} - c_n = (n+1)^2 - 10(n+1) - n^2 + 10n = 2n - 9$, negativo para
$n \leq 4$ e positivo para $n \geq 5$: $(c_n)$ decresce até
$c_5 = -25$, seu \omterm{def:g10:functions:extrema}{mínimo}, e depois cresce. Não é \omterm{def:g12:seq:monotonic}{monótona}.
\end{solution}

\begin{solution}{exo:g12:seq:3}
Coloque os termos dominantes em evidência:
\[
u_n = \frac{n^2(3 - 1/n + 1/n^2)}{n^2(2 + 5/n^2)} \xrightarrow[n\to+\infty]{} \frac{3}{2}.
\]
Multiplique pelo conjugado:
\[
v_n = \frac{(n+1) - n}{\sqrt{n+1} + \sqrt{n}} = \frac{1}{\sqrt{n+1}+\sqrt{n}}
\xrightarrow[n\to+\infty]{} 0.
\]
Divida numerador e denominador por $3^n$:
\[
w_n = \frac{(2/3)^n - 1}{1 + (1/3)^n} \xrightarrow[n\to+\infty]{} \frac{0-1}{1+0} = -1,
\]
usando $\lim q^n = 0$ para $\abs{q} < 1$.
\end{solution}

\begin{solution}{exo:g12:seq:4}
$u_n = 5 + 3n \to +\infty$ e $v_n = 8 \cdot (1/2)^n = 2^{3-n} \to 0$. As
somas são
\[
\sum_{k=0}^{n} u_k = (n+1)\,\frac{5 + (5+3n)}{2} = \frac{(n+1)(10+3n)}{2}
\xrightarrow[n\to+\infty]{} +\infty,
\]
\[
\sum_{k=0}^{n} v_k = 8\,\frac{1 - (1/2)^{n+1}}{1 - 1/2}
= 16\left(1 - \left(\tfrac{1}{2}\right)^{n+1}\right)
\xrightarrow[n\to+\infty]{} 16 .
\]
\end{solution}

\begin{solution}{exo:g12:seq:5}
Como $-1 \leq \cos n \leq 1$,
\[
\frac{n-1}{n+1} \leq \frac{n + \cos n}{n+1} \leq 1,
\]
e $\frac{n-1}{n+1} \to 1$, de modo que o limite é $1$ pelo teorema do
confronto.

Para o segundo limite, escreva
\[
0 \leq \frac{n!}{n^n}
= \frac{1}{n}\cdot\frac{2}{n}\cdots\frac{n}{n}
\leq \frac{1}{n},
\]
porque todo fator $\frac{k}{n}$ com $2 \leq k \leq n$ é no máximo $1$.
Como $\frac1n \to 0$, o teorema do confronto dá
$\frac{n!}{n^n} \to 0$. (A cota geométrica sugerida também funciona: cada
fator com $k \leq n/2$ é no máximo $\frac12$, o que dá a cota mais forte
$(1/2)^{\floor{n/2}}$.)
\end{solution}

\begin{solution}{exo:g12:seq:6}
\emph{1.} $u_0 = 0 \in \intcc{0}{2}$. Se $0 \leq u_n \leq 2$, então
$2 \leq u_n + 2 \leq 4$, logo $\sqrt{2} \leq u_{n+1} \leq 2$; em
particular, $0 \leq u_{n+1} \leq 2$. Por indução, a propriedade vale para
todo $n$.

\emph{2.} $u_{n+1} - u_n = \sqrt{u_n + 2} - u_n$. Para
$x \in \intcc{0}{2}$,
$\sqrt{x+2} \geq x \iff x + 2 \geq x^2 \iff (2-x)(x+1) \geq 0$, o que é
verdadeiro. Logo $(u_n)$ é \omterm{def:g12:seq:monotonic}{crescente}.

\emph{3.} \omterm{def:g12:seq:monotonic}{Crescente} e \omterm{def:g12:seq:bounded}{limitada superiormente} por $2$, $(u_n)$ \omterm{def:g12:seq:limit}{converge}
para algum $\ell \in \intcc{0}{2}$. Passando ao limite em
$u_{n+1} = \sqrt{u_n + 2}$ (a aplicação $x \mapsto \sqrt{x+2}$ é
contínua), vem $\ell = \sqrt{\ell + 2}$, logo $\ell^2 - \ell - 2 = 0$,
isto é, $\ell \in \{-1, 2\}$. Como $\ell \geq 0$, $\lim u_n = 2$.
\end{solution}

\begin{solution}{exo:g12:seq:7}
\emph{1.} Entre duas doses, $40\%$ do medicamento é eliminado, de modo que
a quantidade $u_n$ vira $0.6\,u_n$; a dose seguinte acrescenta $1$
unidade: $u_{n+1} = 0.6\,u_n + 1$.

\emph{2.} $v_{n+1} = u_{n+1} - 2.5 = 0.6\,u_n + 1 - 2.5 = 0.6(u_n - 2.5)
= 0.6\,v_n$: $(v_n)$ é uma \omterm{def:g12:seq:arith-geom}{progressão geométrica} de razão $0.6$ e primeiro
termo $v_0 = 1 - 2.5 = -1.5$. Logo $v_n = -1.5 \times 0.6^n$ e
\[
u_n = 2.5 - 1.5 \times 0.6^n .
\]

\emph{3.} Como $0.6^n \to 0$, $u_n \to 2.5$: a quantidade de medicamento
se estabiliza em $2.5$ unidades.
\end{solution}

\begin{solution}{exo:g12:seq:8}
\emph{1.} $u_0 = 3 > 1$. Se $u_n > 1$, então $u_n + 2 > 0$ e
\[
u_{n+1} - 1 = \frac{4u_n - 1 - u_n - 2}{u_n + 2} = \frac{3(u_n - 1)}{u_n + 2} > 0 .
\]
Por indução, $u_n > 1$ para todo $n$ (e, em particular,
$u_n + 2 \neq 0$, de modo que a \omterm{def:g12:seq:sequence}{sequência} está bem definida).

\emph{2.} Usando a identidade acima,
\[
v_{n+1} = \frac{1}{u_{n+1} - 1} = \frac{u_n + 2}{3(u_n - 1)}
= \frac{(u_n - 1) + 3}{3(u_n - 1)} = \frac{1}{3} + v_n .
\]
Logo $(v_n)$ é uma \omterm{def:g12:seq:arith-geom}{progressão aritmética} de razão $\frac13$ com
$v_0 = \frac{1}{u_0 - 1} = \frac12$.

\emph{3.} $v_n = \frac12 + \frac{n}{3}$, portanto
$u_n = 1 + \frac{1}{v_n} = 1 + \frac{6}{3 + 2n}$. Como
$v_n \to +\infty$, $u_n \to 1$.
\end{solution}

\begin{solution}{exo:g12:seq:9}
\emph{1.} $H_{2n} - H_n = \sum_{k=n+1}^{2n} \frac{1}{k}$ é uma soma de
$n$ termos, cada um pelo menos $\frac{1}{2n}$; logo
$H_{2n} - H_n \geq n \cdot \frac{1}{2n} = \frac12$.

\emph{2.} Por indução em $k$: $H_{2^0} = H_1 = 1 \geq 1$. Se
$H_{2^k} \geq 1 + \frac{k}{2}$, então, aplicando o item 1 com $n = 2^k$,
\[
H_{2^{k+1}} \geq H_{2^k} + \frac12 \geq 1 + \frac{k+1}{2}.
\]
A \omterm{def:g12:seq:sequence}{sequência} $(H_n)$ é \omterm{def:g12:seq:monotonic}{crescente} (cada passo soma $\frac{1}{n+1} > 0$) e a
subsequência $H_{2^k}$ é ilimitada, de modo que $(H_n)$ não é \omterm{def:g12:seq:bounded}{limitada
superiormente}. \omterm{def:g12:seq:monotonic}{Crescente} e ilimitada, ela tende a $+\infty$
(\cref{thm:g12:seq:monotone}).
\end{solution}

\begin{solution}{exo:g12:seq:10}
\emph{1.} A \omterm{def:g12:seq:sequence}{sequência} $d_n = b_n - a_n$ satisfaz
$d_{n+1} - d_n = (b_{n+1} - b_n) - (a_{n+1} - a_n) \leq 0$, de modo que
$(d_n)$ é \omterm{def:g10:functions:variations}{decrescente}; como $d_n \to 0$, obtemos $d_n \geq 0$ para todo
$n$ (uma \omterm{def:g12:seq:sequence}{sequência} \omterm{def:g10:functions:variations}{decrescente} com um termo negativo permaneceria abaixo
dele para sempre, impedindo o limite $0$). Logo $a_n \leq b_n$.

\emph{2.} De $a_n \leq b_n \leq b_0$, a \omterm{def:g12:seq:sequence}{sequência} \omterm{def:g12:seq:monotonic}{crescente} $(a_n)$ é
\omterm{def:g12:seq:bounded}{limitada superiormente} e, portanto, \omterm{def:g12:seq:limit}{converge} para algum $\ell$. Do mesmo
modo, $(b_n)$, \omterm{def:g10:functions:variations}{decrescente} e \omterm{def:g12:seq:bounded}{limitada inferiormente} por $a_0$, \omterm{def:g12:seq:limit}{converge}
para algum $\ell'$. Então $\ell' - \ell = \lim (b_n - a_n) = 0$, logo
$\ell = \ell'$.

\emph{3.} $(a_n)$ é (estritamente) \omterm{def:g12:seq:monotonic}{crescente}, pois
$a_{n+1} - a_n = \frac{1}{(n+1)!} > 0$. Para $(b_n)$,
\[
b_{n+1} - b_n = \frac{1}{(n+1)!} + \frac{1}{(n+1)(n+1)!} - \frac{1}{n\,n!}
= \frac{n(n+1) + n - (n+1)^2}{n(n+1)(n+1)!}
= \frac{-1}{n(n+1)(n+1)!} < 0 ,
\]
de modo que $(b_n)$ é \omterm{def:g10:functions:variations}{decrescente}. Por fim,
$b_n - a_n = \frac{1}{n\,n!} \to 0$. As duas \omterm{def:g12:seq:sequence}{sequências} são adjacentes e,
portanto, convergem para um limite comum.
\end{solution}

\begin{solution}{pb:g12:seq:1}
\textbf{1.} Verdadeiro para $n = 1$ ($1 = 1^2$). Se
$1 + 3 + \dots + (2n - 1) = n^2$, então, somando o ímpar
seguinte: $n^2 + (2n + 1) = (n + 1)^2$: hereditariedade. Por
indução, vale para todo $n \geq 1$.

\textbf{2.} $2^0 = 1 > 0$. Se $2^n > n$, então
$2^{n+1} = 2 \cdot 2^n > 2n \geq n + 1$ para $n \geq 1$ (e
$n = 0$ se verifica diretamente): hereditariedade, pronto.

\textbf{3.} $n = 0$: $1 \geq 1$. Se $(1 + x)^n \geq 1 + nx$,
multiplique por $1 + x \geq 1 > 0$:
$(1 + x)^{n+1} \geq (1 + nx)(1 + x) = 1 + (n + 1)x + nx^2
\geq 1 + (n + 1)x$.

\textbf{4.} O passo de $n = 1$ para $n = 2$: entre duas
bolinhas, ``as $n$ primeiras'' e ``as $n$ últimas'' são duas
bolinhas isoladas e \emph{disjuntas} --- nenhuma bolinha comum
faz a ponte entre os dois grupos, de modo que nada obriga as
cores a coincidirem. O argumento de hereditariedade exige em
silêncio que os dois grupos se sobreponham, o que só é verdade
a partir de $n \geq 2$; com o caso base $n = 1$, a corrente
nunca começa.

\textbf{5.} $4^0 - 1 = 0 = 3 \times 0$. Se $4^n - 1 = 3k$,
então $4^{n+1} - 1 = 4(4^n - 1) + 3 = 3(4k + 1)$:
hereditariedade.

\textbf{6.} $x_1 = \frac32$, $x_2 = \frac{17}{12}$,
$x_3 = \frac{577}{408}$.

\textbf{7.} $x_{n+1}^2 - 2 = \frac{(x_n^2 + 2)^2 - 8x_n^2}
{4x_n^2} = \frac{(x_n^2 - 2)^2}{4 x_n^2}$: um quadrado dividido
por um positivo, logo $\geq 0$, e $> 0$ sempre que
$x_n^2 \neq 2$. Indução: $x_0 = 2 > 0$ com $x_0^2 = 4 > 2$; se
$x_n > 0$ e $x_n^2 > 2$, então $x_{n+1}$ (uma \omterm{def:g11:stat:mean}{média} de
positivos) é positivo e $x_{n+1}^2 - 2 > 0$.

\textbf{8.} $x_{n+1} - x_n = \frac{2 - x_n^2}{2x_n} < 0$ pela
questão 7: estritamente \omterm{def:g10:functions:variations}{decrescente}.

\textbf{9.} \omterm{def:g10:functions:variations}{Decrescente} e \omterm{def:g12:seq:bounded}{limitada inferiormente} (por $1$, já
que $x_n^2 > 2 > 1$ e $x_n > 0$): pelo teorema da convergência
\omterm{def:g12:seq:monotonic}{monótona}, $(x_n)$ \omterm{def:g12:seq:limit}{converge} para algum $L \geq 1$.

\textbf{10.} Os limites respeitam a álgebra: de
$x_{n+1} = \frac12\left(x_n + \frac{2}{x_n}\right)$ e
$x_n \to L \geq 1 > 0$ vem
$L = \frac12\left(L + \frac2L\right)$, logo $L^2 = 2$ e, sendo
$L$ positivo, $L = \sqrt2$. Veredicto: convergência
demonstrada, limite identificado --- Heron absolvido com
honras.

\textbf{11.} $x_{n+1} - \sqrt2 = \frac{x_n^2 - 2\sqrt2\,x_n +
2}{2x_n} = \frac{(x_n - \sqrt2)^2}{2x_n}$: exatamente
$e_{n+1} = \frac{e_n^2}{2x_n}$, e $x_n > \sqrt2$ dá
$e_{n+1} \leq \frac{e_n^2}{2\sqrt2}$. Erro elevado ao quadrado:
cada passo dobra o número de casas decimais corretas, como se
observava desde o ano 9.

\textbf{12.} $e_0 \approx 0.5858$, $e_1 \approx 0.0858$,
$e_2 \approx 0.00245$, $e_3 \approx 2.1 \times 10^{-6}$.
Razões: $\frac{e_1}{e_0^2} \approx 0.25 = \frac{1}{2x_0}$;
$\frac{e_2}{e_1^2} \approx 0.333 = \frac{1}{2x_1}$;
$\frac{e_3}{e_2^2} \approx 0.353 \approx \frac{1}{2x_2}$: o
teorema em ação.

\textbf{13.} $H_{2^{18}} \geq 1 + 9 = 10$: cerca de
$260\,000$ termos ($2^{18} = 262\,144$) apenas para passar de
$10$ --- divergência a passo de tartaruga (e $H_n > 100$
exigiria mais termos do que átomos em qualquer biblioteca).

\textbf{14.} $S_n = 2 - \frac{1}{2^n}$ (soma geométrica), e
$\frac{1}{2^n} \to 0$ (\cref{thm:g12:seq:geometric}), logo
$S_n \to 2$: a barra de chocolate mordida sem fim tende ao
inteiro sem nunca alcançá-lo --- agora na linguagem oficial dos
limites.

\textbf{15.} Para $k \geq 2$:
$\frac{1}{k^2} \leq \frac{1}{k(k-1)} = \frac{1}{k-1} -
\frac1k$, de modo que
$S_n \leq 1 + \left(1 - \frac1n\right) < 2$: \omterm{def:g12:seq:monotonic}{crescente} e
\omterm{def:g12:seq:bounded}{limitada superiormente}, logo convergente (convergência
\omterm{def:g12:seq:monotonic}{monótona}). Euler mais tarde nomeou o limite:
$\frac{\pi^2}{6}$.

\textbf{16.} Os termos tenderem a $0$ é \emph{necessário} para
que as somas se acomodem, mas não decide nada: os termos
harmônicos $\frac1n \to 0$ e, ainda assim, as somas explodem;
os termos $\frac{1}{n^2} \to 0$ e as somas convergem. A questão
toda é o quão \emph{depressa} os termos morrem --- a teoria das
séries, construída nos volumes de graduação.

\textbf{17.} $a_1 = \sqrt2 \approx 1.41421$, $b_1 = 1.5$;
$a_2 \approx 1.45648$, $b_2 \approx 1.45711$: duas iterações já
concordam em três casas decimais --- velocidade impressionante.

\textbf{18.} $b_{n+1} - a_{n+1} = \frac{a_n + b_n}{2} -
\sqrt{a_n b_n} = \frac{(\sqrt{b_n} - \sqrt{a_n})^2}{2} \geq
0$: as \omterm{def:g11:stat:mean}{médias} mantêm a ordem. $(a_n)$ cresce:
$a_{n+1} = \sqrt{a_n b_n} \geq \sqrt{a_n \cdot a_n} = a_n$; e
$(b_n)$ decresce de modo simétrico.

\textbf{19.} $\frac{b_{n+1} - a_{n+1}}{b_n - a_n} =
\frac{(\sqrt{b_n} - \sqrt{a_n})^2}
{2(\sqrt{b_n} - \sqrt{a_n})(\sqrt{b_n} + \sqrt{a_n})}
= \frac{\sqrt{b_n} - \sqrt{a_n}}{2(\sqrt{b_n} + \sqrt{a_n})}
\leq \frac12$: a diferença pelo menos se reduz à metade, logo
$b_n - a_n \to 0$; com a questão 18, as \omterm{def:g12:seq:sequence}{sequências} são
adjacentes e partilham um limite $M(1, 2)$.

\textbf{20.} A terceira iteração dá
$a_3 \approx b_3 \approx 1.456791$: $M(1, 2) \approx
1.456791$ em três voltas de manivela (a diferença praticamente
\emph{se eleva ao quadrado}, como em Heron). Ranking de
velocidade das \omterm{def:g12:seq:sequence}{sequências} deste problema, da mais lenta à mais
rápida: as somas harmônicas (divergência glacial), as somas
geométricas (erro reduzido à metade a cada passo), Heron e a
\omterm{def:g11:stat:mean}{média} aritmético-geométrica (erro elevado ao quadrado a cada
passo) --- e foi a velocidade sobrenatural dessa \omterm{def:g11:stat:mean}{média} que disse
a Gauss que ele havia encontrado um novo veio da análise.
\end{solution}
